

To keep the prerequisite minimal, some basic definitions and notations are listed in this section. For
details and proofs of certain results, see Lickorish's book \cite{LBook} and \cite{BHMV}. For a
survey of the skein theory and some generalization,
see \cite{LSurvey}.

Given a compact oriented 3-manifold $M$, and a ring of coefficients
$R=Z[A,A^{-1}]$, define a $R-module$ by
\[SS(M)=the\ free\ module\ generated\ by\ framed\ links\ in\ M\]
where two framed links isotopic to each other are regarded as the same.

\begin{definition} The skein module associated to $M$ is
\[S(M)= SS(M)/the\ Kauffman\ skein\ relations\]
\end{definition}

These skein relations are the following:

\[\scalebox{.15}{\psfig{file=skeinrelationcircle.eps}}\]


\[\scalebox{.15}{\psfig{file=skeinrelation.eps}}\]


\subsection{Examples}\label{examples}

\subsubsection{The skein module of $S^3$}\label{three sphere}
The skein module of $S^3$ is equal to $Z[A,A^{-1}]$. The reason is the following.

Given a planar diagram of a famed link $L$ in $S^3$, by applying the skein relation repeatedly to resolve the crossings and remove the unknots, $L$ can be reduced to the empty link multiplied by a Laurent polynomial of $A$. In this way, a polynomial is associated to each framed link in $S^3$.

Multiplied by a normalization factor, the dependence of the polynomial on the framing can be removed. Up to change of variables, this is the Jones polynomial of $L$.

\subsubsection{The skein module of $D\times S^1$}\label{solid torus}
The case of the solid torus is slightly different. By simply using the skein
relations, not every link can be reduced to the empty one. Certain
essential links remain.

$S(D\times S^1)$ is found to be a free  $R-module$
generated by those elements represented by parallel strings along
the core of the solid torus, see the pictures below.
\[\scalebox{.35}{\psfig{file=basis1.eps}}\]
\subsubsection{The algebra structure of $S(D\times S^1)$ and an alternative basis}\label{algebra and new basis}
The skein module of the standard solid torus $D\times S^1$ has an
algebra structure. The multiplication is defined as follows.

Fix a curve on $\partial(D\times S^1)$, for instance $\{1\}\times S^1$. Here $D$ is identified with the unit disk in the complex plane. Name this curve $C_+$.
Choose another curve on $\partial(D\times S^1)$ parallel to $C_+$, for instance
$\{-1\}\times S^1$, name it  $C_-$. Given two copies of $D\times
S^1$, glue the first one to the second one by identifying a neighborhood
of $C_+$ in the boundary from the first copy to a neighborhood of
$C_-$ in the boundary from the second copy. The result is still a solid
torus.

Two elements from $S(D\times S^1)$ which can be represented by skeins are multiplied by juxtaposition of the two skeins representing them. See the picture below.
\[\scalebox{.45}{\psfig{file=multiplication.eps}} \]

The multiplication is then extended linearly. With this
multiplication, $S(D\times S^1)$ is a commutative $R-algebra$.
It is isomorphic to $R[z_1]$, where $z_1$ is one of the basis
elements described above. Note that $z_n=z_1^n$.

 Define elements $\{e_n|n\ integers\}$ of $S(D\times S^1)$ by $e_0=1,e_1=z_1$ and $z_1e_n=e_{n-1}+e_{n+1}$.
These relations determine $e_n$ for $n\ge 0$ as well as those $e_n$ where $n<0$. It can be checked by induction that $e_n=-e_{-n-2}$.

The set $\{e_n|n\ge 0 \}$ is a basis of
$S(D\times S^1)$, since each $e_n$ is a monic polynomial of $z_1$ with degree n.

It is worth mentioning the description of the multiplication using the new basis.

\begin{definition} A triple of nonnegative integers $(l,m,n)$ is said to be admissible if $l+m+n$ is even and $(l,m,n)$ satisfies the triangle inequality, that is \[ |m-n|\le l \le m+n\]
\end{definition}

The multiplication of the basis elements is given by
\[ e_m e_n=\sum_{(l,m,n)\ admissible} e_l\]
for $m,n$ nonnegative.
This can also be checked by induction on n.

\subsection{The multilinear bracket}\label{multilinear bracket}
Given an framed link $L$ in $M$ with components $K_1,K_2,...K_m$, denote the framing of $K_j$ by $\sigma_j$. A multilinear
bracket $\langle *,*,...,*\rangle_L$ in $S(M)$ can be constructed as following:

\

The framing $\sigma_j$ of $K_j$ determines a diffeomorphism
\[\tau_j:D\times S^1 \to a\ tubular\ neighbourhood\ of\ K_j.\]
$\tau_j$ gives a map $D\times S^1 \to M$. It induces a homomorphism
\[(\tau_j)_*:S(D\times S^1) \to S(M).\]
These $(\tau_j)_*$ combine
together to give a multi-linear map \[\tau_*: S(D\times
S^1)\times S(D\times S^1)\times ...\times S(D\times S^1)\to S(M)\]
where the product contains $m$ factors.

Define $\langle
x_1,x_2,...,x_m\rangle_L$ to be $\tau_*(x_1,x_2,...,x_m)$.


\

Components of $L$ can be colored by integers. Assume $K_j$ is colored by $n_j$, then the colored framed link $L$ gives an element in $S(M)$, that is $\langle e_{n_1},e_{n_2},...,e_{n_m}\rangle_L$.

\begin{remark}
Change the framing of certain components of $L$, while keeping the colors unchanged, for example $\sigma_1$ is changed to $\sigma_1+b$, where $b$ is an integer. Then
\[\langle e_{n_1},e_{n_2},...,e_{n_m}\rangle_{L'}=(-1)^{bn_1}A^{b(n_1^2+2n_1)}\langle e_{n_1},e_{n_2},...,e_{n_m}\rangle_L\]
where $L'$ is the same link as $L$ with changed framings.
\end{remark}



\subsubsection{Examples}\label{examples}
 Let $K$ be a knot in $S^3$ with framing $\sigma$, the n-th colored Jones polynomial of K, $J_n^{\sigma}(A,K)$, is defined to be $\langle e_n
\rangle_K$,which is an element of $S(S^3)$. Since $S(S^3)$ is equal to $Z[A,A^{-1}]$, this element is a Laurent polynomial in A.

 In particular,
\[\langle e_n\rangle_U=(-1)^n[n+1]\]
where $U$ is the unknot with framing zero, and $[m]=\frac{A^{2m}-A^{-2m}}{A^2-A^{-2}}$.

Let L be the following link in $D\times S^1$:
\[\scalebox{.25}{\psfig{file=link-L1.eps}}\]

Here both components have framing 0, --- consult the framing convention below. Then we have \[\langle e_m,e_n\rangle_L=(-1)^m \frac{[(m+1)(n+1)]}{[n+1]} e_n\]

If $K$ has framing 0, $J_n(A,K)$ is used for the colored Jones polynomial instead of $J_n^0(A,K)$.

\subsubsection{The framing convention}
Since the linking number of two oriented knots in $S^3$ is defined, the framing of a knot in $S^3$ can be specified by an integer,  the linking number.

For an arbitrary 3-manifold $M$, the framing of a knot $K$ in $M$ can't be specified by an integer.
It is only possible under certain circumstances.

For example, $M$
is the solid torus $D\times S^1$ with a chosen curve $C_+$,
$\{1\}\times S^1$, on its boundary. Embed $M$ into $S^3$. A
framed knot $K$ in $M$  is then mapped to a framed knot in $S^3$. Its
framing corresponds to an integer. In general the corresponding
integer will change if the embedding is changed. If $D\times S^1$ is embedded into  $S^3$ in such a way that $C_+$ has linking
number zero with the core of the solid torus after the embedding, the corresponding integer will be the same for all such embeddings.

In the following, a
knot $K$ in $D\times S^1$ is said to have framing n for some integer n in the
sense that the knot is viewed as a knot in $S^3$ through such kind
of embedding of  $D\times S^1$ into $S^3$.







